\section{Conclusion and Open Problems}
\label{sec:conclusion}

Quotient-matrix reductions and exact computational sweeps constrain the
landscape of Leontovich graphs sharply. The impossibility at $n=5$ for
the spherically symmetric family $T(x,y,z)$
(Proposition~\ref{prop:n5}) and the exact checks at $n=6$ identify
$n=7$ as the smallest crossover threshold witnessed here.

We conclude with the two principal open directions highlighted by this work:

\begin{enumerate}
  \item \textbf{Near-path sufficiency.} Every known Leontovich graph is
    detected by the near-path family $E_n^{(d)}$. The near-path
    sufficiency conjecture asks whether this always holds. If it does,
    the bounded sweeps in this paper upgrade directly to unconditional
    nonexistence proofs for small targets.
  \item \textbf{Minimal simple targets.} The bipartite double cover of
    $\hat{T}(1,35,1,50)$ yields a simple strongly Leontovich graph on
    3,644 vertices, and the perturbed quotient
    (Section~\ref{sec:doublecover}) gives high-precision numerical evidence
    for a simple non-bipartite candidate on 6,806 vertices. Closing the
    gap to the 12-vertex frontier probed by the simple-graph sweep
    remains a computational challenge.
\end{enumerate}
